\subsection{პროექტის ტექნიკური მხარე}

\subsubsection{სისტემის საერთო არქიტექტურა}

Peared-ის სისტემა სამი გაშვებადი (runtime) კომპონენტისგან შედგება:

\begin{enumerate}[itemsep=0pt]
  \item \textbf{desktop აპლიკაცია} (\texttt{tauri-app}), ეს ნაწილი ეშვება
        ყველა მონაწილის მიერ. სწორედ აქ ცხოვრობს დოკუმენტის CRDT ასლი, Monaco
        რედაქტორი და სესიის მართვის მთელი ლოგიკა.
  \item \textbf{გეითვეი} (\texttt{gateway}), Go-ზე დაწერილი WebSocket
        relay სერვერი, რომელსაც \emph{მხოლოდ ჰოსტის} აპლიკაცია უშვებს
        დამხმარე (sidecar) პროცესად სესიის დაწყებისას და კლავს სესიის
        დასრულებისას. ის მდგომარეობას თითქმის არ ინახავს, უბრალოდ
        გადააგზავნის ბაიტებს ოთახის წევრებს შორის.
  \item \textbf{Cloudflared Tunnel}, არასავალდებულო დამხმარე პროცესი,
        რომელიც ჰოსტის გეითვეის დროებით საჯარო
        \texttt{*.trycloudflare.com} მისამართზე აქვეყნებს, რათა ლოკალური
        ქსელის გარეთ მყოფმა სტუმრებმაც შეძლონ შეერთება.
\end{enumerate}

ქსელური ტოპოლოგია hub-and-spoke მოდელს მიჰყვება: ყველა მონაწილე (თავად
ჰოსტის კლიენტიც კი) ერთსა და იმავე გეითვეის უკავშირდება WebSocket-ით, ხოლო
გეითვეი მიღებულ ყოველ ფრეიმს დანარჩენებს გადაუგზავნის. ამით ვიღებთ
peer-to-peer სისტემის თვისებებს (დოკუმენტი მხოლოდ მონაწილეებთანაა,
ცენტრალური სერვისი არ არსებობს) სრული mesh-ქსელის სირთულის გარეშე -
თითოეულ კლიენტს მხოლოდ ერთი კავშირი სჭირდება.

\begin{diagram}[htbp]
  \centering
  \resizebox{\textwidth}{!}{%
  \begin{tikzpicture}[
      font=\small\sffamily,
      box/.style={draw, thick, rounded corners=3pt, align=center,
                  minimum width=3.2cm, minimum height=1.15cm},
      hostapp/.style={box, draw=pcblue, fill=pcblueL},
      side/.style={box, draw=pcgreen, fill=pcgreenL},
      tun/.style={box, draw=pcorange, fill=pcorangeL},
      guest/.style={box, draw=pcpurple, fill=pcpurpleL},
      cloud/.style={box, draw=pcorange, fill=pcorangeL!50, dashed},
      arr/.style={{Stealth}-{Stealth}, very thick, draw=black!70},
      lbl/.style={font=\footnotesize\sffamily, fill=white, inner sep=1.5pt}
    ]
    % host machine
    \node[hostapp] (app) {Peared app\\(host)};
    \node[side, right=2.6cm of app] (gw) {Gateway\\(Go relay)};
    \node[tun, below=1.0cm of gw] (cf) {cloudflared};

    % guests
    \node[guest, right=3.4cm of gw] (g1) {Peared app\\(LAN guest)};
    \node[cloud, below=1.9cm of g1] (edge) {Cloudflare edge\\ *.trycloudflare.com};
    \node[guest, right=1.3cm of edge] (g2) {Peared app\\(remote guest)};

    \begin{scope}[on background layer]
      \node[draw=pcblue!60, thick, dashed, rounded corners=5pt, fill=pcblue!6,
            inner sep=10pt, fit=(app)(gw)(cf),
            label={[font=\small\sffamily\bfseries, text=pcblue]above:Host machine}] {};
    \end{scope}

    \draw[arr, draw=pcblue] (app) -- node[lbl]{ws://localhost} (gw);
    \draw[arr, draw=pcpurple] (gw) -- node[lbl]{LAN WebSocket} (g1);
    \draw[arr, draw=pcorange] (cf.east) -- node[lbl, sloped]{tunnel} (edge.west);
    \draw[arr, draw=pcpurple] (edge) -- node[lbl]{WSS} (g2);
    \draw[arr, dashed, draw=pcorange] (gw) -- (cf);
  \end{tikzpicture}}
  \caption{სისტემის საერთო არქიტექტურა: ჰოსტის კომპიუტერზე გაშვებული
  გეითვეი და მასთან დაკავშირებული კლიენტები}
  \label{diag:architecture}
\end{diagram}

\paragraph{სესიის სასიცოცხლო ციკლი.}
როდესაც მომხმარებელი სესიას იწყებს (hosting), აპლიკაცია
თანმიმდევრულად ასრულებს შემდეგ ნაბიჯებს: (1) უშვებს გეითვეის პროცესს და
მისგან stdout-ზე დაბეჭდილი JSON შეტყობინებით იგებს,
რომელ პორტზე დაიწყო მუშაობა; (2) უშვებს
\texttt{cloudflared}-ს და ელოდება საჯარო ბმულს; (3) HTTP მოთხოვნით
გეითვეისგან ქმნის „ოთახს“ (room) და იღებს მის იდენტიფიკატორს; (4) თავად
უერთდება ოთახს WebSocket-ით და მაშინვე აგზავნის დოკუმენტის სნეფშოტს, რათა
მოგვიანებით შემოსული სტუმრები სწრაფად „დაეწიონ“ მიმდინარე მდგომარეობას.
ამის შემდეგ ჰოსტი ორ ბმულს იღებს,- ლოკალური ქსელის და საჯაროს.
სტუმრის მხარეს იგივე პროცესი ბევრად
მარტივია: ბმულის ჩასმა, WebSocket კავშირი, სნეფშოტის მიღება და მუშაობის
დაწყება.

სესიაში უფლებები ასიმეტრიულია: ჰოსტი ყოველთვის სრულუფლებიანია, ხოლო
სტუმრები, თუ ჰოსტმა არ მიუთითა, მხოლოდ კითხვის რეჟიმში ერთვებიან (ჰოსტს შეუძლია
სესიის დაწყებისას აირჩიოს, რომ სტუმრებს ჩაწერის უფლება ავტომატურად
მიეცეთ). სესიის მიმდინარეობისას ჰოსტი თითოეულ სტუმარს ინდივიდუალურად
ანიჭებს ან ართმევს ჩაწერის უფლებას, ამ მექანიზმის ტექნიკურ დეტალებს
ქვემოთ განვიხილავთ.

\subsubsection{Tauri ფრეიმვორკი კონცეპტუალურად}

desktop აპლიკაცია აგებულია Tauri 2 ფრეიმვორკზე \cite{tauri}. Tauri-ის
იდეა მარტივია: აპლიკაციის ინტერფეისი იწერება ჩვეულებრივი ვებ-ტექნოლოგიებით
(ჩვენს შემთხვევაში React + TypeScript) და ეშვება ოპერაციული სისტემის
ჩაშენებულ webview-ში, ხოლო „ნატიური“ ლოგიკა, ფაილებთან წვდომა, ქსელი,
პროცესების მართვა, იწერება Rust-ზე და ცალკე პროცესის ბირთვში ცხოვრობს.
ორ სამყაროს შორის კომუნიკაცია ხდება IPC-ით: ფრონტენდი იძახებს ბექენდის
„ბრძანებებს“ (commands), ბექენდი კი ფრონტენდს „ივენთებს“ (events) უგზავნის.
Electron-თან შედარებით Tauri-ის ბინარები გაცილებით მცირეა (ის Chromium-ს
არ აპაკეტებს) და ბექენდის ენად Rust გამოვიყენეთ.
CRDT ბიბლიოთეკა სწორედ Rust-ზეა დაწერილი და აპლიკაციაში პირდაპირ,
დანამატების გარეშე ჩაშენდა.

Tauri-ის კიდევ ერთი ფუნქცია, რომელსაც აქტიურად ვიყენებთ, არის
\emph{sidecar} ბინარები: გეითვეი და \texttt{cloudflared} აპლიკაციასთან
ერთად იდება და აპლიკაცია მათ სასიცოცხლო ციკლს სრულად აკონტროლებს -
მომხმარებელი ვერც კი ამჩნევს, რომ მის კომპიუტერზე ცალკე სერვერი გაეშვა.

\subsubsection{ფრონტენდი და Monaco, როგორც ბაზისი}

ინტერფეისის ბაზისად Monaco Editor \cite{monaco} ავირჩიეთ, იგივე
 კომპონენტი, რომელზეც VS Code-ია აშენებული. ეს არჩევანი ბევრ რამეს
გვაძლევს „უფასოდ“: სინტაქსის გაფერადება, მრავალკურსორიანი რედაქტირება,
undo/redo, ძებნა და keyboard shortcuts. ჩვენი ფრონტენდი
Monaco-ს გარშემოა აგებული: React მართავს აპლიკაციის შელს (გვერდითი
პანელი, სესიის სტატუსი, მონაწილეთა სია, ფაილების მენიუ), ხოლო Monaco -
თავად ტექსტს. ლოკალური ცვლილებები Monaco-დან ბექენდისკენ IPC ბრძანებებით
მიდის, remote მონაწილეების ცვლილებები კი ბექენდიდან ივენთებით ბრუნდება.

\begin{diagram}[htbp]
  \centering
  \begin{tikzpicture}[
      font=\small\sffamily,
      layer/.style={draw, thick, rounded corners=3pt, align=center,
                    minimum height=1.05cm, minimum width=3.9cm},
      arr/.style={{Stealth}-{Stealth}, very thick, draw=black!70}
    ]
    % Frontend
    \node[layer, draw=pcblue, fill=pcblueL, minimum width=13.0cm]
      (fe) {\textbf{Frontend:} React + Monaco Editor (webview)};
    \node[below=0.5cm of fe, font=\small\itshape] (ipc)
      {IPC: commands $\downarrow$ \quad events $\uparrow$};

    % Backend core row 1
    \node[layer, draw=pcblue, fill=pcblueL, below=0.4cm of ipc, xshift=-4.55cm]
      (appstate) {AppState\\session FSM};
    \node[layer, draw=pcgreen, fill=pcgreenL, right=0.45cm of appstate]
      (actor) {Document actor\\(crdt-core + tokio)};
    \node[layer, draw=pcpurple, fill=pcpurpleL, right=0.45cm of actor]
      (ws) {WebSocket\\management};

    % Backend row 2
    \node[layer, draw=pcblue, fill=pcblueL, below=0.55cm of appstate]
      (perm) {Roster \&\\permissions};
    \node[layer, draw=pcorange, fill=pcorangeL, right=0.45cm of perm]
      (gc) {Garbage collection\\(SV exchange)};
    \node[layer, draw=pcgreen, fill=pcgreenL, right=0.45cm of gc]
      (pers) {Persistence\\(file open/save)};

    % Backend row 3
    \node[layer, draw=pcpurple, fill=pcpurpleL, below=0.55cm of perm]
      (proc) {Process\\coordinator};
    \node[layer, draw=pcorange, fill=pcorangeL, right=0.45cm of proc]
      (metrics) {Metrics\\aggregator};
    \node[layer, draw=pcblue, fill=pcblueL, right=0.45cm of metrics]
      (gwapi) {Gateway API\\client (HTTP)};

    \begin{scope}[on background layer]
      \node[draw=black!50, thick, dashed, rounded corners=5pt, fill=black!3,
            inner sep=9pt, fit=(appstate)(ws)(proc)(gwapi),
            label={[font=\small\sffamily\bfseries]below:Tauri backend (Rust)}] {};
    \end{scope}

    \draw[arr, draw=pcblue] (fe) -- (ipc);
    \draw[arr, draw=pcblue] (ipc) -- (actor);
  \end{tikzpicture}
  \caption{კლიენტის (Tauri აპლიკაციის) შრეები და ქვესისტემები}
  \label{diag:layers}
\end{diagram}
